% Preámbulo común para las soluciones de los prácticos de Programación 3
\documentclass[11pt,a4paper]{article}

\usepackage[utf8]{inputenc}
\usepackage[T1]{fontenc}
\usepackage{lmodern}
\usepackage[spanish,es-noshorthands]{babel}
\usepackage[margin=2.3cm]{geometry}
\usepackage{amsmath,amssymb,amsthm}
\usepackage{enumitem}
\usepackage{xcolor}
\usepackage{tikz}
\usetikzlibrary{arrows.meta,positioning,shapes,calc}
\usepackage[noend]{algpseudocode}
\usepackage{algorithm}
\usepackage{booktabs}
\usepackage{tcolorbox}
\usepackage[hidelinks]{hyperref}

\floatname{algorithm}{Algoritmo}
\renewcommand{\algorithmicrequire}{\textbf{Entrada:}}
\renewcommand{\algorithmicensure}{\textbf{Salida:}}
\algrenewcommand\algorithmicwhile{\textbf{mientras}}
\algrenewcommand\algorithmicdo{\textbf{hacer}}
\algrenewcommand\algorithmicfor{\textbf{para}}
\algrenewcommand\algorithmicforall{\textbf{para todo}}
\algrenewcommand\algorithmicif{\textbf{si}}
\algrenewcommand\algorithmicthen{\textbf{entonces}}
\algrenewcommand\algorithmicelse{\textbf{sino}}
\algrenewcommand\algorithmicreturn{\textbf{retornar}}
\algrenewcommand\algorithmicfunction{\textbf{función}}
\algrenewcommand\algorithmicprocedure{\textbf{procedimiento}}

\theoremstyle{plain}
\newtheorem{lema}{Lema}[section]
\newtheorem{teorema}[lema]{Teorema}
\newtheorem{proposicion}[lema]{Proposición}
\theoremstyle{definition}
\newtheorem{definicion}[lema]{Definición}
\theoremstyle{remark}
\newtheorem*{observacion}{Observación}

% Caja para el enunciado de cada ejercicio
\newtcolorbox{enunciado}{colback=gray!6,colframe=gray!50,boxrule=0.4pt,arc=2pt,left=6pt,right=6pt,top=4pt,bottom=4pt,fontupper=\small}
% Nota que vincula el ejercicio con parciales/exámenes
\newcommand{\examref}[1]{\par\noindent{\small\color{blue!50!black}\textbf{Relación con parciales/exámenes:} #1}\par\medskip}

\newcommand{\NP}{\mathsf{NP}}
\newcommand{\Pclass}{\mathsf{P}}
\newcommand{\Oh}{\mathcal{O}}

\setlength{\parskip}{0.4em}
\setlength{\parindent}{0pt}

\setlength{\emergencystretch}{2em}

\newcommand{\IS}{\textsc{Independent Set}}
\newcommand{\VC}{\textsc{Vertex Cover}}
\newcommand{\SP}{\textsc{Set Packing}}
\newcommand{\PSP}{\textsc{PSP}}
\newcommand{\IPROG}{\textsc{(0--1) IPROG}}
\newcommand{\SATMON}{\textsc{SAT-MON}}
\newcommand{\SAT}{\textsc{SAT}}
\newcommand{\leP}{\le_P}
\newcommand{\SI}{S\'I}

\title{Práctico 12 --- Problemas NP-Completos}
\author{Programación 3 -- FING, UdelaR}
\date{}

\begin{document}
\maketitle

\section*{Esquema de prueba}

Un problema de decisión $X$ es \emph{NP-completo} si (1) $X\in\NP$ y (2) para todo $Y\in\NP$, $Y\leP X$.
Como $\leP$ es transitiva, para (2) alcanza con elegir un problema $Z$ que ya sabemos NP-completo y probar
$Z\leP X$ (¡en ese sentido: \emph{desde} el conocido \emph{hacia} el nuevo!). Siguiendo el formato de las
soluciones oficiales de los parciales:

\begin{enumerate}
  \item \textbf{$X\in\NP$.} Definir un certificado $t$ y un certificador $B(s,t)$ y mostrar que: $B$ corre en
  tiempo polinomial en $|s|+|t|$; para instancias \SI{} existe un certificado de largo polinomial en $|s|$ con el
  que $B$ responde \SI{}; y si $B(s,t)$ responde \SI{} para algún $t$, entonces $s$ es \SI.
  \item \textbf{$Z\leP X$.} Describir una transformación de una instancia genérica $s$ de $Z$ en una instancia
  $f(s)$ de $X$, mostrar que es polinomial, y probar las dos direcciones: $s$ \SI{} $\Rightarrow$ $f(s)$ \SI, y
  $f(s)$ \SI{} $\Rightarrow$ $s$ \SI.
  \item \textbf{Concluir} por transitividad.
\end{enumerate}

%======================================================================
\section{Ejercicio 1}

\begin{enunciado}
Consideramos un sistema \textbf{lineal} de inecuaciones sobre un conjunto de variables
$\{x_1,x_2,\dots,x_n\}$. El \emph{Problema de la Programación Entera Binaria, (0--1) IPROG}, consiste en
decidir si existe una asignación de valores enteros en $\{0,1\}$ a las variables tal que el sistema se
satisface. Por ejemplo, para el sistema
\[
\begin{cases}
2x_1+x_2+x_3 \le 3\\
x_1 \le 1\\
x_1-2x_2+x_3 \ge 2
\end{cases}
\]
la respuesta al problema de decisión es SÍ, porque la asignación $x_1=1$, $x_2=0$, $x_3=1$ satisface todas
las inecuaciones.
\begin{enumerate}[label=(\alph*)]
  \item Demuestre que (0--1) IPROG pertenece a la clase $\NP$.
  \item Demuestre que $\mathit{SAT}\leP$ (0--1) IPROG, donde recordamos que SAT es el problema de
  \emph{Satisfacción Booleana}.

  \textbf{Sugerencia:} Para construir una reducción de tiempo polinomial
  \begin{itemize}
    \item Construya una instancia de (0--1) IPROG con la misma cantidad de variables que la de SAT.
    \item Interprete los valores de verdad como enteros del conjunto $\{0,1\}$.
    \item Exprese cada cláusula de la instancia de SAT como una inecuación entre variables de la instancia de (0--1) IPROG.
  \end{itemize}
  \item Concluya que (0--1) IPROG es NP-Completo.
\end{enumerate}
\end{enunciado}

\examref{No aparece literalmente. El formato de la prueba de pertenencia a $\NP$ (certificado = tira de
valores para las variables, evaluación de cada restricción) sigue al del Segundo parcial 2024, Ej.~2 (4-SAT) y
del Segundo parcial 2020, Ej.~3 (casiSAT); la reducción desde SAT es análoga a la de Examen febrero 2022,
Ej.~2 (FALSABLE, $\SAT\leP\text{FALSABLE}$).}

\paragraph{Representación.} Una instancia de \IPROG{} es un sistema de $r$ inecuaciones
\[
  \sum_{j=1}^n a_{ij}\,x_j \ \bowtie_i\ b_i, \qquad 1\le i\le r,\quad \bowtie_i\in\{\le,\ge\},
\]
dado por la matriz entera $A=(a_{ij})$ ($r\times n$), el vector $b$ y los $r$ símbolos $\bowtie_i$. Su tamaño es
$N=\Theta(rn)$ números, cada uno con a lo sumo $L$ bits.

\subsection*{(a) \IPROG{} $\in\NP$}

\paragraph{Certificado.} Una tira $t=t_1t_2\cdots t_n$ de $n$ bits, donde $t_j$ es el valor asignado a $x_j$.
Su largo, $n$, es menor o igual que el tamaño de la instancia (la matriz tiene $n$ columnas), así que es
polinomial en $|s|$.

\paragraph{Certificador.} Recibe la instancia y el certificado y evalúa cada inecuación con los valores
dados por el certificado:
\begin{algorithm}[H]
\caption{Certificador de \IPROG}
\begin{algorithmic}[1]
\Require Instancia $s=(A,b,\bowtie)$ con $r$ inecuaciones y $n$ variables; certificado $t\in\{0,1\}^n$
\If{$|t|\neq n$} \State \Return NO \EndIf
\For{$i\gets 1$ \textbf{a} $r$}
  \State $S\gets 0$
  \For{$j\gets 1$ \textbf{a} $n$}
    \If{$t_j=1$} \State $S\gets S+a_{ij}$ \EndIf
  \EndFor
  \If{no se cumple $S\bowtie_i b_i$} \State \Return NO \EndIf
\EndFor
\State \Return SÍ
\end{algorithmic}
\end{algorithm}

\paragraph{Tiempo.} Se hacen $rn$ sumas. Cada suma parcial es la suma de a lo sumo $n$ enteros de $L$ bits,
así que tiene $\Oh(L+\log n)$ bits y cada suma o comparación cuesta $\Oh(L+\log n)$. El tiempo total es
$\Oh(rn(L+\log n))$, polinomial en $|s|+|t|$.

\paragraph{Correctitud.}
\begin{itemize}
  \item Si $s$ es \SI{}, existe una asignación $x\in\{0,1\}^n$ que satisface todas las inecuaciones. Con el
  certificado $t=x$, para cada $i$ el valor $S$ calculado es $\sum_j a_{ij}x_j$, que cumple $S\bowtie_i b_i$;
  el certificador no retorna NO en ninguna iteración y responde SÍ.
  \item Si existe $t$ con el que el certificador responde SÍ, entonces $|t|=n$ y para cada $i$ se cumplió
  $\sum_j a_{ij}t_j\bowtie_i b_i$. Luego la asignación $x_j=t_j\in\{0,1\}$ satisface el sistema y $s$ es \SI.
\end{itemize}
Por lo tanto existe un certificador eficiente y $\IPROG\in\NP$.

\subsection*{(b) $\SAT\leP\IPROG$}

Sea $(X,\mathcal C)$ una instancia de \SAT{} con variables $X=\{x_1,\dots,x_n\}$ y cláusulas
$\mathcal C=\{C_1,\dots,C_k\}$. Para cada cláusula $C_i$ sean
$P_i$ la lista (multiconjunto) de índices $j$ tales que $x_j$ aparece como término positivo en $C_i$, y $N_i$ la
lista de índices $j$ tales que $\overline{x_j}$ aparece en $C_i$.

\paragraph{Transformación.} Construimos un sistema sobre las mismas $n$ variables $x_1,\dots,x_n$
(interpretando $1=$ verdadero, $0=$ falso), con una inecuación por cláusula:
\[
  (I_i)\qquad \sum_{j\in P_i} x_j \;+\; \sum_{j\in N_i} (1-x_j)\ \ge\ 1,
\]
que, pasando las constantes al lado derecho para que quede en la forma lineal estándar, es
\[
  (I_i)\qquad \sum_{j\in P_i} x_j \;-\; \sum_{j\in N_i} x_j\ \ge\ 1-|N_i|.
\]
(Si una variable aparece varias veces se suman los coeficientes; por ejemplo
$x_1\vee\overline{x_2}\vee\overline{x_3}$ se traduce en $x_1-x_2-x_3\ge -1$.)

\paragraph{Tiempo polinomial.} Se recorre cada cláusula una vez y, por cada término, se suma $+1$ o $-1$ a un
coeficiente y se cuenta $|N_i|$. Con la matriz $k\times n$ de coeficientes (inicializada en 0), el tiempo es
$\Oh(kn+\ell)$, donde $\ell$ es la cantidad total de términos. Los coeficientes tienen valor absoluto a lo sumo
$\ell$. Todo es polinomial en el tamaño de la instancia de \SAT.

\paragraph{Lema clave.} Sea $\sigma$ una asignación de verdad y $x^\sigma\in\{0,1\}^n$ el vector con
$x^\sigma_j=1\iff\sigma(x_j)=$ verdadero. Para cada término: un término positivo $x_j$ es verdadero si y solo si
$x^\sigma_j=1$, y un término negativo $\overline{x_j}$ es verdadero si y solo si $1-x^\sigma_j=1$. Como cada
sumando vale 0 o 1, el lado izquierdo de la primera forma de $(I_i)$ evaluado en $x^\sigma$ es
\emph{exactamente la cantidad de términos verdaderos de $C_i$} (contados con multiplicidad). Por lo tanto
\[
  \sigma \text{ satisface } C_i \iff \text{hay al menos un término verdadero en } C_i \iff x^\sigma \text{ satisface } (I_i).
\]

\paragraph{Equivalencia.} La correspondencia $\sigma\mapsto x^\sigma$ es una biyección entre asignaciones de
verdad y vectores de $\{0,1\}^n$.
\begin{enumerate}[label=(\roman*)]
  \item Si $(X,\mathcal C)$ es \SI{}, existe $\sigma$ que satisface todas las cláusulas; por el lema $x^\sigma$
  satisface todas las inecuaciones $(I_1),\dots,(I_k)$ y el sistema es \SI.
  \item Si el sistema es \SI{}, existe $x\in\{0,1\}^n$ que satisface todas las $(I_i)$; sea $\sigma$ la
  asignación con $x^\sigma=x$. Por el lema, $\sigma$ satisface cada $C_i$, y $(X,\mathcal C)$ es \SI.
\end{enumerate}
Así \SAT{} se resuelve construyendo el sistema en tiempo polinomial y llamando una vez a la caja negra de
\IPROG: $\SAT\leP\IPROG$.

\paragraph{Ejemplo.} La instancia $(x_1\vee x_2)\wedge(\overline{x_1}\vee x_3)\wedge(\overline{x_2}\vee\overline{x_3})$
se transforma en
\[
  x_1+x_2\ge1,\qquad -x_1+x_3\ge 0,\qquad -x_2-x_3\ge -1,
\]
y la asignación $x_1=1,x_2=0,x_3=1$ satisface ambas instancias.

\subsection*{(c) \IPROG{} es NP-completo}

Por (a), $\IPROG\in\NP$. Por el teorema de Cook--Levin, \SAT{} es NP-completo, es decir, $Y\leP\SAT$ para todo
$Y\in\NP$. Por (b), $\SAT\leP\IPROG$, y por transitividad de $\leP$, $Y\leP\IPROG$ para todo $Y\in\NP$. Ambas
condiciones juntas dicen que \IPROG{} es NP-completo.

%======================================================================
\section{Ejercicio 2}

\begin{enunciado}
(Kleinberg \& Tardos, Ex.~8.9.) Consideramos una red de comunicaciones, modelada por un grafo dirigido
$G=(V,E)$. Hay $c$ usuarios a los cuales les interesa usar la red. Cada usuario $i$, $1\le i\le c$, hace un
requerimiento para reservar un camino específico $P_i$ en $G$ sobre el cual transmitir sus datos. Al operador
de la red le interesa aceptar la máxima cantidad de solicitudes de caminos posible, bajo la restricción de
que no puede haber nodos en común entre ningún par de caminos que sean aceptados.

Definimos el \emph{Path Selection Problem} (PSP) de la siguiente manera. Dado un grafo dirigido $G=(V,E)$, un
conjunto de caminos $\{P_1,P_2,\dots,P_c\}$ en $G$ y un entero $k$, $0\le k\le c$, ¿es posible seleccionar al
menos $k$ de dichos caminos de modo que ningún par de ellos tenga nodos en común? Muestre que PSP es
NP-completo.

\textbf{Sugerencia:} Busque una reducción desde el problema \emph{Set Packing}.
\end{enunciado}

\examref{Muy parecido al Segundo parcial 2025, Ej.~2 (UnionEsV: ¿hay $k$ caminos cuya unión de vértices es
$V$?), donde la solución oficial reduce desde \emph{Set Cover} con la misma idea: un vértice por elemento del
universo y un camino por subconjunto. Aquí la condición es de \emph{disjunción} en lugar de
\emph{cubrimiento}, por eso se reduce desde \SP.}

\paragraph{Problema de referencia.} \SP{} (K\&T, Sección~8.2): dado un conjunto $U=\{u_1,\dots,u_n\}$, una
colección $S_1,\dots,S_\ell$ de subconjuntos de $U$ y un entero $k$, ¿existen al menos $k$ de esos conjuntos
que sean disjuntos dos a dos? \SP{} es NP-completo: está en $\NP$ y $\IS\leP\SP$ (K\&T, Cap.~8) tomando
$U=E$ y, por cada vértice $v$, el conjunto $S_v$ de aristas incidentes a $v$ (dos vértices son adyacentes si y
solo si sus conjuntos se intersecan).

\subsection*{\PSP{} $\in\NP$}

Una instancia $s$ es el grafo $G=(V,E)$, los caminos $P_1,\dots,P_c$ (cada uno como lista de nodos) y $k$.

\paragraph{Certificado.} Una tira $t$ de $c$ bits, $t_i=1$ si y solo si el camino $P_i$ es seleccionado.
Su largo $c$ es a lo sumo el tamaño de la instancia.

\paragraph{Certificador.} Recibe la instancia y el certificado, verifica la cantidad de caminos elegidos y
marca los nodos usados por los caminos elegidos, rechazando si un nodo se usa dos veces:
\begin{algorithm}[H]
\caption{Certificador de \PSP}
\begin{algorithmic}[1]
\Require $s=(G,\{P_1,\dots,P_c\},k)$, $t\in\{0,1\}^c$
\If{$|t|\ne c$ \textbf{o} la cantidad de unos de $t$ es menor que $k$} \State \Return NO \EndIf
\State $\mathit{usado}[v]\gets\mathit{false}$ para todo $v\in V$
\For{$i\gets1$ \textbf{a} $c$ \textbf{tal que} $t_i=1$}
  \ForAll{nodo $v$ de $P_i$}
     \If{$\mathit{usado}[v]$} \State \Return NO \EndIf
     \State $\mathit{usado}[v]\gets\mathit{true}$
  \EndFor
\EndFor
\State \Return SÍ
\end{algorithmic}
\end{algorithm}
(Suponemos, como es usual, que los caminos son simples, es decir, que no repiten nodos.)

\paragraph{Tiempo.} $\Oh(|V|+c+\sum_i|P_i|)$, lineal en el tamaño de la entrada.

\paragraph{Correctitud.} Si $s$ es \SI{}, sea $J$ un conjunto de al menos $k$ índices de caminos disjuntos dos
a dos en nodos; con $t_i=1\iff i\in J$ el certificador pasa el chequeo de cardinal y nunca encuentra un nodo ya
usado (eso pasaría solo si dos caminos de $J$ compartieran un nodo), así que responde SÍ. Recíprocamente, si
responde SÍ con $t$, el conjunto $J=\{i:t_i=1\}$ tiene al menos $k$ elementos y ningún nodo fue marcado dos
veces, es decir, los caminos de $J$ son disjuntos en nodos dos a dos: $s$ es \SI. Luego $\PSP\in\NP$.

\subsection*{$\SP\leP\PSP$}

\paragraph{Transformación.} Sea $(U,\{S_1,\dots,S_\ell\},k)$ una instancia de \SP{} con $U=\{u_1,\dots,u_n\}$.
\begin{itemize}
  \item Si $k>\ell$ la instancia es trivialmente NO. Como \PSP{} exige $k\le c$, devolvemos una instancia NO
  fija de \PSP: el grafo con nodos $\{a,b\}$ y la arista $(a,b)$, los caminos $P_1=P_2=(a,b)$ y $k=2$ (los
  únicos dos caminos comparten nodos). Si $k\le 0$ ambas instancias son \SI{} y la construcción de abajo sirve.
  \item Si $0\le k\le\ell$ construimos $(G,\{P_1,\dots,P_\ell\},k)$ así:
  \begin{enumerate}
    \item \textbf{Nodos.} Un nodo $w_j$ por cada elemento $u_j\in U$, y dos nodos \emph{privados} $a_i$, $z_i$ por
    cada conjunto $S_i$:
    $V=\{w_1,\dots,w_n\}\cup\{a_1,z_1,\dots,a_\ell,z_\ell\}$.
    \item \textbf{Caminos y aristas.} Si $S_i=\{u_{j_1},\dots,u_{j_r}\}$ con $j_1<j_2<\dots<j_r$, el camino es
    \[
      P_i = a_i \to w_{j_1}\to w_{j_2}\to\cdots\to w_{j_r}\to z_i
    \]
    (si $S_i=\emptyset$, $P_i=a_i\to z_i$), y $E$ es el conjunto de todas las aristas dirigidas usadas por
    algún $P_i$.
    \item $c=\ell$ y el mismo $k$.
  \end{enumerate}
\end{itemize}

Cada $P_i$ es un camino dirigido \emph{simple} de $G$: sus aristas están en $E$ por definición, y no repite
nodos porque los índices $j_1<\dots<j_r$ son estrictamente crecientes y $a_i,z_i$ aparecen una vez. Los nodos
privados garantizan que todo camino tenga al menos una arista aun si $S_i$ es vacío o unitario. Además:
\begin{equation}\label{eq:nodos}
  \text{nodos}(P_i) = \{a_i,z_i\}\cup\{w_j : u_j\in S_i\}. \tag{$*$}
\end{equation}

\begin{center}
\begin{tikzpicture}[>=Stealth, v/.style={circle,draw,inner sep=1pt,minimum size=6mm}, p/.style={rectangle,draw,inner sep=2pt,minimum size=5mm,fill=gray!15}]
  \node[v] (w1) at (0,0) {$w_1$};
  \node[v] (w2) at (2,0) {$w_2$};
  \node[v] (w3) at (4,0) {$w_3$};
  \node[v] (w4) at (6,0) {$w_4$};
  \node[p] (a1) at (-1.5,1.3) {$a_1$}; \node[p] (z1) at (3.5,1.6) {$z_1$};
  \node[p] (a2) at (0.5,-1.6) {$a_2$}; \node[p] (z2) at (7.5,-1.3) {$z_2$};
  \node[p] (a3) at (4.5,1.6) {$a_3$}; \node[p] (z3) at (7.5,1.3) {$z_3$};
  \draw[->,thick,blue] (a1)--(w1); \draw[->,thick,blue] (w1)--(w2); \draw[->,thick,blue] (w2)--(z1);
  \draw[->,thick,red] (a2)--(w2); \draw[->,thick,red] (w2) to[bend right=25] (w4); \draw[->,thick,red] (w4)--(z2);
  \draw[->,thick,green!50!black] (a3)--(w3); \draw[->,thick,green!50!black] (w3) to[bend left=25] (z3);
  \node[align=left,anchor=west] at (8.5,0) {\small$S_1=\{u_1,u_2\}$ \textcolor{blue}{(azul)}\\
    \small$S_2=\{u_2,u_4\}$ \textcolor{red}{(rojo)}\\ \small$S_3=\{u_3\}$ \textcolor{green!50!black}{(verde)}};
\end{tikzpicture}
\end{center}
En el ejemplo, $S_1$ y $S_3$ son disjuntos y también lo son $P_1$ y $P_3$; $S_1\cap S_2=\{u_2\}$ y $P_1$, $P_2$
comparten el nodo $w_2$.

\paragraph{Tiempo polinomial.} $|V|=n+2\ell$; cada camino tiene $|S_i|+2$ nodos y $|S_i|+1$ aristas. Ordenar
los elementos de cada $S_i$ (o recorrerlos con un arreglo de marcas de tamaño $n$) y escribir los caminos lleva
$\Oh(\sum_i |S_i|\log|S_i| + n+\ell)$, y $|E|\le\sum_i(|S_i|+1)$. Todo es polinomial en el tamaño de la
instancia de \SP.

\paragraph{Lema clave.} Para $i\ne i'$: $P_i$ y $P_{i'}$ tienen un nodo en común $\iff$ $S_i\cap S_{i'}\ne\emptyset$.

\emph{Prueba.} Por \eqref{eq:nodos}, y como los nodos privados $a_i,z_i$ no pertenecen a ningún otro camino,
$\text{nodos}(P_i)\cap\text{nodos}(P_{i'}) = \{w_j : u_j\in S_i\cap S_{i'}\}$, que es no vacío si y solo si
$S_i\cap S_{i'}\neq\emptyset$. \qed

\paragraph{Equivalencia} (caso $k\le\ell$; el caso $k>\ell$ es NO en ambos lados por construcción).
\begin{enumerate}[label=(\roman*)]
  \item Si $(U,\{S_i\},k)$ es \SI{} de \SP, existe $J\subseteq\{1,\dots,\ell\}$, $|J|\ge k$, con los $S_i$,
  $i\in J$, disjuntos dos a dos. Por el lema, los caminos $P_i$, $i\in J$, son disjuntos en nodos dos a dos, y son
  al menos $k$: la instancia de \PSP{} es \SI.
  \item Si la instancia de \PSP{} es \SI, existe $J$, $|J|\ge k$, con los $P_i$, $i\in J$, sin nodos en común dos a
  dos. Por el lema, los $S_i$, $i\in J$, son disjuntos dos a dos: la instancia de \SP{} es \SI.
\end{enumerate}
Luego $\SP\leP\PSP$.

\paragraph{Conclusión.} $\PSP\in\NP$, \SP{} es NP-completo y $\SP\leP\PSP$; por transitividad todo problema de
$\NP$ se reduce polinomialmente a \PSP, así que \PSP{} es NP-completo.

\begin{observacion}
También se puede reducir directamente desde \IS{} (componiendo $\IS\leP\SP\leP\PSP$): un nodo $w_e$ por
arista $e$ de $G$ y un camino por vértice $v$ que pasa por los $w_e$ de las aristas incidentes a $v$. Como en la
solución oficial de UnionEsV, en lugar de agregar sólo las aristas usadas se podría tomar el grafo dirigido
completo sobre los $w_j$; lo importante es que cada $P_i$ sea efectivamente un camino de $G$.
\end{observacion}

%======================================================================
\section{Ejercicio 3}

\begin{enunciado}
(Kleinberg \& Tardos, Ex.~8.6.) Consideramos el problema de \emph{Satisfacción Booleana} (SAT),
especificado por un conjunto de cláusulas, $C_1,\dots,C_k$, cada una de las cuales es una disyunción de
términos sobre variables booleanas $x_1,\dots,x_n$. Se dice que una instancia es \emph{monótona} si cada
término de cada cláusula consiste en una variable no negada, esto es, ningún término es de la forma
$\overline{x_i}$.

Las instancias monótonas de SAT son muy fáciles de resolver, ya que asignándole valor de verdad 1 a todas
las variables las cláusulas siempre se satisfacen. Por ejemplo, el conjunto de cláusulas
$\{(x_1\vee x_2),(x_1\vee x_3),(x_2\vee x_3)\}$ define una instancia monótona de SAT, y la asignación de 1 a
todas las variables satisface todas las cláusulas. Sin embargo, observamos que asignar el valor 1 a las
variables $x_1$ y $x_2$ es suficiente para satisfacer las tres cláusulas, independientemente del valor que
asignemos a $x_3$. Resulta entonces natural preguntarse cuál es la cantidad mínima de variables $x_i$ que es
necesario fijar en 1 para satisfacer una determinada instancia monótona de SAT.

Dada una instancia monótona de SAT y un entero $k$, el problema de decisión \emph{Satisfacción Monótona
con Pocas Variables Verdaderas} (SAT-MON) pregunta: ¿hay alguna asignación de valores de verdad en la cual a
lo sumo a $k$ variables se les asigna el valor 1 y todas las cláusulas se satisfacen? Muestre que SAT-MON es
NP-completo.

\textbf{Sugerencia:} Busque una reducción desde un problema sobre grafos.
\end{enunciado}

\examref{Es esencialmente la misma reducción que la del Segundo parcial 2022, Ej.~3 (CONTRATAR, con
sugerencia ``use Vertex Cover''): allí cada arista $(v_j,v_k)$ se transforma en una tarea con aplicantes
calificados $\{a_j,a_k\}$; aquí, en una cláusula $(x_j\vee x_k)$. La prueba de pertenencia a $\NP$ sigue el
formato de 4-SAT (Segundo parcial 2024, Ej.~2).}

\subsection*{\SATMON{} $\in\NP$}

Una instancia $s$ consiste en las variables $x_1,\dots,x_n$, las cláusulas $C_1,\dots,C_k$ (listas de índices
de variables, todas no negadas) y el entero $K$ (lo llamamos $K$ para no confundirlo con la cantidad de
cláusulas).

\paragraph{Certificado.} Una tira $t$ de $n$ bits: $t_j$ es el valor asignado a $x_j$. Su largo es polinomial
(a lo sumo el tamaño de la instancia, suponiendo, como en la solución oficial de 4-SAT, que toda variable aparece
en alguna cláusula o en la lista de variables).

\paragraph{Certificador.} $B(s,t)$: (1) si $|t|\ne n$ o la cantidad de unos de $t$ es mayor que $K$, responde
NO; (2) para cada cláusula $C_i$, recorre sus términos y verifica que alguno $x_j$ tenga $t_j=1$; si alguna
cláusula no tiene ninguno, responde NO; (3) si no, responde SÍ.

\paragraph{Tiempo.} $\Oh(n+\ell)$, donde $\ell$ es la cantidad total de términos: polinomial en $|s|+|t|$.

\paragraph{Correctitud.} Si $s$ es \SI{}, existe una asignación con a lo sumo $K$ variables en 1 que satisface
todas las cláusulas; tomándola como $t$, el paso (1) no rechaza y en cada cláusula hay un término verdadero,
así que $B$ responde SÍ. Si $B(s,t)$ responde SÍ, entonces $t$ tiene a lo sumo $K$ unos y satisface toda
cláusula; es una asignación que prueba que $s$ es \SI. Luego $\SATMON\in\NP$.

\subsection*{$\VC\leP\SATMON$}

\paragraph{Transformación.} Sea $(G,K)$ una instancia de \VC{} con $G=(V,E)$, $V=\{v_1,\dots,v_n\}$,
$E=\{e_1,\dots,e_q\}$. Construimos la instancia monótona:
\begin{enumerate}
  \item una variable $x_j$ por cada vértice $v_j$, $1\le j\le n$;
  \item una cláusula $C_i=(x_a\vee x_b)$ por cada arista $e_i=(v_a,v_b)$, $1\le i\le q$;
  \item el mismo entero $K$.
\end{enumerate}
Todas las cláusulas tienen solo términos no negados: la instancia es monótona.

\paragraph{Tiempo polinomial.} Se crean $n$ variables y se recorre $E$ una vez escribiendo una cláusula de dos
términos por arista: $\Oh(n+q)$ con listas de adyacencia ($\Oh(n^2)$ con matriz). Polinomial.

\paragraph{Correspondencia.} A cada $S\subseteq V$ le asociamos la asignación $\sigma_S$ con
$\sigma_S(x_j)=1\iff v_j\in S$, y viceversa (es una biyección). El número de variables en 1 de $\sigma_S$ es
$|S|$, y
\[
  \sigma_S \text{ satisface } C_i=(x_a\vee x_b) \iff v_a\in S \text{ o } v_b\in S \iff S \text{ cubre } e_i.
\]

\paragraph{Equivalencia.}
\begin{enumerate}[label=(\roman*)]
  \item Si $(G,K)$ es \SI{} de \VC, existe un vertex cover $S$ con $|S|\le K$. Cada arista $e_i=(v_a,v_b)$ tiene
  un extremo en $S$, así que $\sigma_S$ satisface $C_i$; como $i$ es arbitrario, $\sigma_S$ satisface todas las
  cláusulas con $|S|\le K$ variables en 1: la instancia de \SATMON{} es \SI.
  \item Si la instancia de \SATMON{} es \SI{}, existe una asignación $\sigma$ con a lo sumo $K$ variables en 1 que
  satisface todas las cláusulas. Sea $S=\{v_j:\sigma(x_j)=1\}$, de modo que $\sigma=\sigma_S$ y $|S|\le K$. Para
  cada arista $e_i=(v_a,v_b)$, $\sigma$ satisface $C_i=(x_a\vee x_b)$, luego $v_a\in S$ o $v_b\in S$: $S$ es un
  vertex cover de tamaño $\le K$ y $(G,K)$ es \SI.
\end{enumerate}

\paragraph{Ejemplo.} El ejemplo del enunciado es exactamente la imagen del triángulo $v_1v_2v_3$: las cláusulas
$(x_1\vee x_2),(x_1\vee x_3),(x_2\vee x_3)$ corresponden a sus tres aristas, y la asignación
$x_1=x_2=1$ corresponde al vertex cover $\{v_1,v_2\}$ (mínimo, de tamaño 2).

\begin{center}
\begin{tikzpicture}[v/.style={circle,draw,inner sep=1pt,minimum size=6mm}, sel/.style={v,fill=blue!20}]
  \node[sel] (1) at (0,0) {$v_1$};
  \node[sel] (2) at (2,0) {$v_2$};
  \node[v] (3) at (1,1.6) {$v_3$};
  \draw (1)--node[below]{\small$x_1\vee x_2$}(2);
  \draw (1)--node[left]{\small$x_1\vee x_3$}(3);
  \draw (2)--node[right]{\small$x_2\vee x_3$}(3);
\end{tikzpicture}
\end{center}

\paragraph{Conclusión.} $\SATMON\in\NP$, \VC{} es NP-completo y $\VC\leP\SATMON$ (transformación polinomial
más una llamada a la caja negra). Por transitividad, $Y\leP\SATMON$ para todo $Y\in\NP$: \SATMON{} es
NP-completo.

\begin{observacion}
La restricción ``a lo sumo $K$ variables en 1'' es la que da la dificultad: sin ella, una instancia monótona
siempre es satisfacible (todas en 1). Notar también que la reducción va \emph{desde} \VC{} \emph{hacia}
\SATMON; la transformación inversa (cada cláusula como ``arista'') no sirve directamente porque las cláusulas
pueden tener más de dos términos (serían hiperaristas), y además iría en el sentido equivocado.
\end{observacion}

\end{document}
